\section{AHSME 1997}
        \begin{enumerate}
        	\item (\textit{AHSME 1997}) If $\texttt{a}$ and $\texttt{b}$ are digits for which


 $\begin{array}{ccc}& 2 & a\\ \times & b & 3\\ \hline & 6 & 9\\ 9 & 2\\ \hline 9 & 8 & 9\end{array}$


 then $\texttt{a+b =}$


 $\mathrm{(A)\ } 3 \qquad \mathrm{(B) \ }4 \qquad \mathrm{(C) \  } 7 \qquad \mathrm{(D) \  } 9 \qquad \mathrm{(E) \  }12$


 
 \textbf{Answer: }C



\textbf{Solution}

From the units digit calculation, we see that the units digit of $a\times 3$ is $9$.  Since $0 \le a \le 9$ and $a$ is an integer, the only value of $a$ that works is is $a=3$.  As a double-check, that does work, since $23 \times 3 = 69$, which is the first line of the multiplication.

 The second line of the multiplication can be found by doing the multiplication $23\times b = 92$.  Dividing both sides by $23$ gives $b=4$.

 Thus, $a + b = 3 + 4 = 7$, and the answer is $\boxed{C}$.

 


	\item (\textit{AHSME 1997}) The adjacent sides of the decagon shown meet at right angles.  What is its perimeter?


 
                    \begin{figure}[H]
                        centering
                        \includegraphics[width=0.5\textwidth]{Problem/AHSME_I_1997_2_Problem_Figure_1.png}
                    \end{figure}
                    


 $\mathrm{(A)\ } 22 \qquad \mathrm{(B) \ }32 \qquad \mathrm{(C) \  } 34 \qquad \mathrm{(D) \  } 44 \qquad \mathrm{(E) \  }50$


 
 \textbf{Answer: }D



\textbf{Solution}

The three unlabelled vertical sides have the same sum as the two labelled vertical sides, which is $10$.

 The four unlabelled horizontal sides have the same sum as the one large horizontal side, which is $12$.

 Thus, the perimeter is $2(12+10) = 44$, which is option $\boxed{D}$.

 


	\item (\textit{AHSME 1997}) If $x$, $y$, and $z$ are real numbers such that


 $(x-3)^2 + (y-4)^2 + (z-5)^2 = 0$,

 then $x + y + z =$


 $\mathrm{(A)\ } -12 \qquad \mathrm{(B) \ }0 \qquad \mathrm{(C) \  } 8 \qquad \mathrm{(D) \  } 12 \qquad \mathrm{(E) \  }50$


 



 
 \textbf{Answer: }D



\textbf{Solution}

If the sum of three squared expressions is zero, then each expression itself must be zero, since $a^2 \ge 0$ with the equality iff $a=0$.

 In this case, $x-3=0$, $y-4=0$, and $z-5=0$.  Adding the three equations and moving the constant to the right gives $x + y + z = 12$, and the answer is $\boxed{D}$.

 


	\item (\textit{AHSME 1997}) If $a$ is $50\%$ larger than $c$, and $b$ is $25\%$ larger than $c$, then $a$ is what percent larger than $b$?


 $\mathrm{(A)\ } 20\% \qquad \mathrm{(B) \ }25\% \qquad \mathrm{(C) \  } 50\% \qquad \mathrm{(D) \  } 100\% \qquad \mathrm{(E) \  }200\%$


 
 \textbf{Answer: }A



\textbf{Solution 1}

Translating each sentence into an equation, $a = 1.5c$ and $b = 1.25c$.

 We want a relationship between $a$ and $b$.  Dividing the second equation into the first will cancel the $c$, so we try that and get:

 $\frac{a}{b} = \frac{1.5}{1.25}$

 $\frac{a}{b} = \frac{150}{125}$

 $\frac{a}{b} = \frac{6}{5}$

 $a = 1.2b$

 In this case, $a$ is $1.2 - 1 = 0.2 = 20\%$ bigger than $b$, and the answer is $\boxed{A}$.

 



\textbf{Solution 2}

Arbitrarily assign a value to one of the variables.  Since $c$ is the smallest variable, let $c = 100$.

 If $a$ is $50\%$ larger than $c$, then $a = 150$.

 If $b$ is $25\%$ larger than $c$, then $b = 125$.

 We see that $\frac{a}{b} = \frac{150}{125} = 1.2$  So, $a$ is $20\%$ bigger than $b$, and the answer is $\boxed{A}$.

 


	\item (\textit{AHSME 1997}) A rectangle with perimeter $176$ is divided into five congruent rectangles as shown in the diagram.  What is the perimeter of one of the five congruent rectangles?

                    \begin{figure}[H]
                        centering
                        \includegraphics[width=0.5\textwidth]{Problem/AHSME_I_1997_5_Problem_Figure_1.png}
                    \end{figure}
                    


 $\mathrm{(A)\ } 35.2 \qquad \mathrm{(B) \ }76 \qquad \mathrm{(C) \  } 80 \qquad \mathrm{(D) \  } 84 \qquad \mathrm{(E) \  }86$


 



 
 \textbf{Answer: }C



\textbf{Solution}

Let $l$ represent the length of one of the smaller rectangles, and let $w$ represent the width of one of the smaller rectangles, with $w < l$.

 From the large rectangle, we see that the top has length $3w$, the right has length $l + w$, the bottom has length $2l$, and the left has length $l + 2$.

 Since the perimeter of the large rectangle is $176$, we know that $172 = 3w + l + w + 2l + l + w$, or $172 = 5w + 4l$

 From the top and bottom of the large rectangle, we know that $3w = 2l$, or $l = 1.5w$.

 Plugging that into the first equation, we get $176 = 5w + 4(1.5)w$

 $176 = 11w$

 $w = 16$

 $l = 1.5w = 24$

 $P = 2l + 2w = 2(16 + 24) = 80$, and the answer is $\boxed{C}$

 



\textbf{Solution 2 (faster, similar, no long yap)}

Call the length $x$ and width $y$.$x+y$ is the width of the big rectangle, and $3y=2x$ is the length (2 lengths can form the big length, as can 3 widths).Now we have a system of equations:$\left\{\begin{matrix}3y=2x\\2(3y+(x+y))=176\end{matrix}\right.$This, after simplification, yields $80$!

 


	\item (\textit{AHSME 1997}) Consider the sequence 


 $1,-2,3,-4,5,-6,\ldots,$


 whose $n$th term is $(-1)^{n+1}\cdot n$. What is the average of the first $200$ terms of the sequence?


 $\textbf{(A)}-\!1\qquad\textbf{(B)}-\!0.5\qquad\textbf{(C)}\ 0\qquad\textbf{(D)}\ 0.5\qquad\textbf{(E)}\ 1$


 



 
 \textbf{Answer: }B



\textbf{Solution}

The average of a list is the sum of all numbers divided by the size of the list.

 The sum of the list can be found by adding the numbers in pairs:  $(1 + -2) + (3 + -4) + ... + (199 + -200)$

 The sum of each pair is $-1$, and there are $100$ pairs, so the total sum is $-100$.

 There are $200$ numbers on the list, so the average is $\frac{-100}{200} = -0.5$, and the answer is $\boxed{B}$

 


	\item (\textit{AHSME 1997}) The sum of seven integers is $-1$. What is the maximum number of the seven integers that can be larger than $13$?


 $\textbf{(A)}\ 1\qquad\textbf{(B)}\ 4\qquad\textbf{(C)}\ 5\qquad\textbf{(D)}\ 6\qquad\textbf{(E)}\ 7$


 



 
 \textbf{Answer: }D



\textbf{Solution}

If the first six integers are $14$, the last number can be $(-14\cdot 6) - 1 = -85$.  The sum of all seven integers will be $-1$.

 However, if all seven integers are over $13$, the smallest possible sum is $14\cdot 7 = 98$.

 Thus, the answer is $6$, which is option $\boxed{D}$.

 


	\item (\textit{AHSME 1997}) Mientka Publishing Company prices its bestseller Where's Walter? as follows: 


 $C(n) =\left\{\begin{matrix}12n, &\text{if }1\le n\le 24\\ 11n, &\text{if }25\le n\le 48\\ 10n, &\text{if }49\le n\end{matrix}\right.$


 where $n$ is the number of books ordered, and $C(n)$ is the cost in dollars of $n$ books. Notice that $25$ books cost less than $24$ books. For how many values of $n$ is it cheaper to buy more than $n$ books than to buy exactly $n$ books?


 $\textbf{(A)}\ 3\qquad\textbf{(B)}\ 4\qquad\textbf{(C)}\ 5\qquad\textbf{(D)}\ 6\qquad\textbf{(E)}\ 8$


 



 
 \textbf{Answer: }D



\textbf{Solution}

Clearly, the areas of concern are where the piecewise function shifts value.  

 Since $C(25) = 11\cdot 25 = 275$, we want to find the least value of $n$ for which $C(n) > 275$.

 If $n \le 24$, then $C(n) = 12n$, so for $C(n) > 275$, $12n > 275$, which is equivalent to $n > 22.91$.  Thus, both $n=23$ and $n=24$ will be more expensive than $n=25$.

 Since $C(49) = 10\cdot 49 = 490$, we want to find the least value of $n$ for which $C(n) > 490$.

 If $25 \le n \le 48$, then $C(n) = 11n$, so for $C(n) > 490$, we have $11n > 490$, leading to $n > 44.5$.  Thus, $n=45, 46, 47, 48$ will be more expensive than $n=49$.

 Thus, there are $2 + 4 = \boxed{6}$ values of $n$ where it's cheaper to buy more books, and the answer is $\boxed{D}$.

 


	\item (\textit{AHSME 1997}) In the figure, $ABCD$ is a $2 \times 2$ square, $E$ is the midpoint of $\overline{AD}$, and $F$ is on $\overline{BE}$.  If $\overline{CF}$ is perpendicular to $\overline{BE}$, then the area of quadrilateral $CDEF$ is


 
                    \begin{figure}[H]
                        centering
                        \includegraphics[width=0.5\textwidth]{Problem/AHSME_I_1997_9_Problem_Figure_1.png}
                    \end{figure}
                    


 $\textbf{(A)}\ 2\qquad\textbf{(B)}\ 3-\frac{\sqrt{3}}{2}\qquad\textbf{(C)}\ \frac{11}{5}\qquad\textbf{(D)}\ \sqrt{5}\qquad\textbf{(E)}\ \frac{9}{4}$


 
 \textbf{Answer: }C



\textbf{Solution 1}

Since $\angle EBA = \angle FCB$ and $\angle FBC = \angle AEB$, we have $\triangle ABE \sim \triangle FCB$.

 $\frac{AB}{FC} = \frac{BE}{CB} = \frac{EA}{BF}$

 $\frac{2}{FC} = \frac{\sqrt{5}}{2} = \frac{1}{BF}$

 From those two equations, we find that $CF = \frac{4}{\sqrt{5}}$ and $BF = \frac{2}{\sqrt{5}}$

 Now that we have $BF$ and $CF$, we can find the area of the bottom triangle $\triangle CFB$:  $\frac{1}{2}\cdot \frac{4}{\sqrt{5}} \cdot \frac{2}{\sqrt{5}} = \frac{4}{5}$

 The area of left triangle $\triangle BEA$ is $\frac{1}{2}\cdot 2 \cdot 1 = 1$

 The area of the square is $4$.

 Thus, the area of the remaining quadrilateral is $4 - 1 - \frac{4}{5} = \frac{11}{5}$, and the answer is $\boxed{C}$

 



\textbf{Solution 2}

Place the square on a coordinate grid so that $B(0,0)$ and $D(2,2)$.  Line $BE$ is $y = 2x$.  Line $CF$ goes through $(2,0)$ and has slope $-\frac{1}{2}$, so it must be $y = 1 -\frac{x}{2}$ 

 The intersection of the two lines is $F$, and $F$ thus has coordinates $(\frac{2}{5}, \frac{4}{5})$.  The altitude from $F$ to $BC$ thus has length $\frac{4}{5}$, so the area of the triangle $BCF$ is $\frac{1}{2}\cdot 2\cdot \frac{4}{5} = \frac{4}{5}$.

 The other triangle has area $1$, and the whole square has area $4$.  As above, we find the area of the quadrilateral by subtracting the two triangles, and we get $\frac{11}{5}$, which is $\boxed{C}$.

 


	\item (\textit{AHSME 1997}) Two six-sided dice are fair in the sense that each face is equally likely to turn up. However, one of the dice has the $4$ replaced by $3$ and the other die has the $3$ replaced by $4$ . When these dice are rolled, what is the probability that the sum is an odd number?


 $\textbf{(A)}\ \frac{1}{3}\qquad\textbf{(B)}\ \frac{4}{9}\qquad\textbf{(C)}\ \frac{1}{2}\qquad\textbf{(D)}\ \frac{5}{9}\qquad\textbf{(E)}\ \frac{11}{18}$


 
 \textbf{Answer: }D



\textbf{Solution}

On the first die, the chance of an odd number is $\frac{4}{6} = \frac{2}{3}$, and the chance of an even number is $\frac{1}{3}$.

 On the second die, the chance of an odd number is $\frac{1}{3}$, and the chance of an even number is $\frac{2}{3}$.

 To get an odd sum, we need exactly one even and one odd.

 The odds of the first die being even and the second die being odd is $\frac{1}{3}\cdot\frac{1}{3} = \frac{1}{9}$

 The odds of the first die being odd and the second die being even is $\frac{2}{3}\cdot\frac{2}{3} = \frac{4}{9}$.

 The overall probability is $\frac{1}{9} + \frac{4}{9} = \frac{5}{9}$, and the answer is $\boxed{D}$.

 


	\item (\textit{AHSME 1997}) In the sixth, seventh, eighth, and ninth basketball games of the season, a player scored $23$,$14$, $11$, and $20$ points, respectively. Her points-per-game average was higher after nine games than it was after the first five games. If her average after ten games was greater than $18$, what is the least number of points she could have scored in the tenth game?


 $\textbf{(A)}\ 26\qquad\textbf{(B)}\ 27\qquad\textbf{(C)}\ 28\qquad\textbf{(D)}\ 29\qquad\textbf{(E)}\ 30$


 
 \textbf{Answer: }D



\textbf{Solution}

The sum of the scores for games $6$ through $9$ is $68$.  The average in these four games is $\frac{68}{4} = 17$.

 The total points in all ten games is greater than $10\cdot 18 = 180$.  Thus, it must be at least $181$.

 There are at least $181 - 68 = 113$ points in the other six games:  games $1-5$ and game $10$.

 Games $1-5$ must have an average of less than $17$.  Thus we cannot put more than $16 + 17 + 17 + 17 + 17 = 84$ points in those five games.

 Thus, the tenth game must have at least $113 - 84 = 29$ points, and the answer is $\boxed{D}$.

 


	\item (\textit{AHSME 1997}) If $m$ and $b$ are real numbers and $mb>0$, then the line whose equation is $y=mx+b$ \textit{cannot} contain the point


 $\textbf{(A)}\ (0,1997)\qquad\textbf{(B)}\ (0,-1997)\qquad\textbf{(C)}\ (19,97)\qquad\textbf{(D)}\ (19,-97)\qquad\textbf{(E)}\ (1997,0)$


 
 \textbf{Answer: }E



\textbf{Solution 1}

Geometrically, $b$ is the y-intercept, and $m$ is the slope.

 Since $mb>0$, then either we have a positive y-intercept and a positive slope, or a negative y-intercept and a negative slope.

 Lines with a positive y-intercept and positive slope can go through quadrants I, II, and III.  They cannot go through quadrant IV, because if you start at $(0,b)$ for a positve $b$, you can't go down into the fourth quadrant with a positive slope.  Thus, point $C$ is a possible point.

 Lines with a negative y-intercept and negative slope can go through quadrants II, III, and IV.  Thus, point $D$ is a possible point.

 Looking at the axes, any point on the y-axis is possible.  Thus, $A$ and $B$ are both possible.

 However, points on the positive x-axis are inpossible to reach.  If you start with a positive y-intercept, you must go up and to the right.  If you start with a negative y-intercept, you must go down and to the right.  Thus, $\boxed{E}$ cannot be reached.

 



\textbf{Solution 2}

Algebraically, if $mb > 0$, then either $m>0$ and $b>0$, or $m<0$ and $b<0$.

 Constructing lines that hit some of the points:

 $y = x + 1997$ will hit $(0,1997)$, so $A$ is possible.

 $y = -x - 1997$ will hit $(0,-1997)$, so $B$ is possible.

 $y = x + 78$ will hit $(19,97)$, so $C$ is possible.

 $y = -x - 78$ will hit $(19, -97)$, so $D$ is possible.

 By process of elimination, $\boxed{E}$ must be impossible.  Plugging in the point $(1997,0)$ into $y = mx + b$ will give $0 = 1997m + b$, or $b = -1997m$.  Thus, $m$ and $b$ must be of opposite signs.

 



\textbf{Solution 3}

Plugging in $(0,1997)$ into $y = mx + b$ gives $1997 = 0m + b$.  Thus, $b$ is positive, and $m$ can be anything, so $m$ can be positive too.

 Plugging in $(0, -1997)$ into $y = mx + b$ gives $-1997 = 0m + b$.  Thus, $b$ is negative, and $m$ can be anything, so $m$ can be negative too.

 Plugging in $(19,97)$ into $y = mx + b$ gives $97 = 19m + b$.  Examining this shows $m$ and $b$ can both be positive.

 Plugging in $(19,-97)$ into $y = mx + b$ gives $-97 = 19m + b$.  Examining this shows $m$ and $b$ can both be negative.

 Plugging in $(1997,0)$ into $y = mx + b$ gives $0 = 1997m + b$.  Examining this shows exactly one of $m$ or $b$ can be positive - the other must be negative.  Thus, the answer is $\boxed{E}$.

 


	\item (\textit{AHSME 1997}) How many two-digit positive integers $N$ have the property that the sum of $N$ and the number obtained by reversing the order of the digits of $N$ is a perfect square?


 $\textbf{(A)}\ 4\qquad\textbf{(B)}\ 5\qquad\textbf{(C)}\ 6\qquad\textbf{(D)}\ 7\qquad\textbf{(E)}\ 8$


 
 \textbf{Answer: }E



\textbf{Solution}

Let $N = 10t + u$, where $t$ is the tens digit and $u$ is the units digit.

 The condition of the problem is that $10t + u + 10u + t$ is a perfect square.

 Simplifying and factoring, we want $11(t+u)$ to be a perfect square.

 Thus, $t+u$ must at least be a multiple of $11$, and since $t$ and $u$ are digits, the only multiple of $11$ that works is $11$ itself.

 Thus, $(t,u) = (2,9)$ is the first solution, and $(t,u) = (9,2)$ is the last solution.  There are $8$ solutions in total, leading to answer $\boxed{E}$.

 


	\item (\textit{AHSME 1997}) The number of geese in a flock increases so that the difference between the populations in year $n+2$ and year $n$ is directly proportional to the population in year $n+1$. If the populations in the years $1994$, $1995$, and $1997$ were $39$, $60$, and $123$, respectively, then the population in $1996$ was


 $\textbf{(A)}\ 81\qquad\textbf{(B)}\ 84\qquad\textbf{(C)}\ 87\qquad\textbf{(D)}\ 90\qquad\textbf{(E)}\ 102$


 
 \textbf{Answer: }B



\textbf{Solution}

Let $x$ be the population in $1996$, and let $k$ be the constant of proportionality.

 If $n=1994$, then the difference in population between $1996$ and $1994$ is directly proportional to the population in $1995$.

 Translating this sentence, $(x - 39) = k(60)$

 Similarly, letting $n=1995$ gives the sentence $(123 - 60) = kx$

 Since $kx = 63$, we have $k = \frac{63}{x}$

 Plugging this into the first equation, we have:

 $(x - 39) = \frac{60\cdot 63}{x}$

 $x - 39 = \frac{3780}{x}$

 $x^2 - 39x - 3780 = 0$

 $(x - 84)(x + 45) = 0$

 Since $x>0$, we must have $x=84$, and the answer is $\boxed{B}$.

 


	\item (\textit{AHSME 1997}) Medians $BD$ and $CE$ of triangle $ABC$ are perpendicular, $BD=8$, and $CE=12$.  The area of triangle $ABC$ is 


 
                    \begin{figure}[H]
                        centering
                        \includegraphics[width=0.5\textwidth]{Problem/AHSME_I_1997_15_Problem_Figure_1.png}
                    \end{figure}
                    


 $\textbf{(A)}\ 24\qquad\textbf{(B)}\ 32\qquad\textbf{(C)}\ 48\qquad\textbf{(D)}\ 64\qquad\textbf{(E)}\ 96$


 
 \textbf{Answer: }D



\textbf{Solution 1}


                    \begin{figure}[H]
                        centering
                        \includegraphics[width=0.5\textwidth]{Solution/AHSME_I_1997_15_Solution_Figure_1.png}
                    \end{figure}
                    

 One median divides a triangle into $2$ equal areas, so all three medians will divide a triangle into $6$ equal areas.

 The median $CE$ is divided into a $2:1$ ratio at centroid $G$, so $GE = \frac{1}{3}\cdot CE = \frac{1}{3}\cdot 12 = 4$

 Similarly, $BG = \frac{2}{3}\cdot 8 = \frac{16}{3}$

 The area of the right triangle $\triangle BEG$ is $\frac{1}{2}\cdot\frac{16}{3}\cdot 4$

 The area of the whole figure is $6\cdot \frac{1}{2}\cdot\frac{16}{3}\cdot 4 = 64$, and the correct answer is $\boxed{D}$.

 



\textbf{Solution 2}

Notice that if you were to draw in line ED, you would get an orthodiagonal quadrilateral with diagonals 8 and 12. The area is going to be equal to 48. Now we need to examine the triangle AED. If the area we are trying to find is denoted as A, we can tell that the area of AEC is A/2. The area of AED is going to be half of that since AD = DC so it would be A/4. This means that 48 is 3/4 of A, so naturally A is going to be 64. Giving $\boxed{D}$

 


	\item (\textit{AHSME 1997}) The three row sums and the three column sums of the array 


 
 \[\left[\begin{matrix}4 & 9 & 2\\ 8 & 1 & 6\\ 3 & 5 & 7\end{matrix}\right]\]


 are the same. What is the least number of entries that must be altered to make all six sums different from one another?


 $\textbf{(A)}\ 1\qquad\textbf{(B)}\ 2\qquad\textbf{(C)}\ 3\qquad\textbf{(D)}\ 4\qquad\textbf{(E)}\ 5$


 
 \textbf{Answer: }D



\textbf{Solution}

If you change $3$ numbers, then you either change one number in each column and row (ie sudoku-style):

 
 \[\left[\begin{matrix}* & 9 & 2\\ 8 & * & 6\\ 3 & 5 & *\end{matrix}\right]\]

 Or you leave at least one row and one column unchanged:

 
 \[\left[\begin{matrix}* & 9 & 2\\ * & * & 6\\ 3 & 5 & 7\end{matrix}\right]\]

 In the first case, you are changing just one common number in two sums, so you wind up with three pairs of sums.  (In the example given, the sum in row $x$ is the same as in column $x$.)

 In the second case, since two of the sums are unchanged, and the sums started out equal, they must remain equal.  (In the second example given, row $3$ and column $3$ are untouched.)

 Either way, $3$ changes is not enough. However, building on the second example, if you change either the untouched column or the untouched row, you will get a possible answer:

 
 \[\left[\begin{matrix}* & 9 & 2\\ * & * & 6\\ 3 & * & 7\end{matrix}\right]\]

 Letting the $*$ be a zero does indeed give $6$ different sums, so the answer is $4$, which is option $\boxed{D}$.

 


	\item (\textit{AHSME 1997}) A line $x=k$ intersects the graph of $y=\log_5 x$ and the graph of $y=\log_5 (x + 4)$. The distance between the points of intersection is $0.5$. Given that $k = a + \sqrt{b}$, where $a$ and $b$ are integers, what is $a+b$?


 $\textbf{(A)}\ 6\qquad\textbf{(B)}\ 7\qquad\textbf{(C)}\ 8\qquad\textbf{(D)}\ 9\qquad\textbf{(E)}\ 10$


 
 \textbf{Answer: }A



\textbf{Solution}

Since the line $x=k$ is vertical, we are only concerned with vertical distance.

 In other words, we want to find the value of $k$ for which the distance $|\log_5 x - \log_5 (x+4)| = \frac{1}{2}$

 Since $\log_5 x$ is a strictly increasing function, we have:

 $\log_5 (x + 4) - \log_5 x = \frac{1}{2}$

 $\log_5 (\frac{x+4}{x}) = \frac{1}{2}$

 $\frac{x+4}{x} = 5^\frac{1}{2}$

 $x + 4 = x\sqrt{5}$

 $x\sqrt{5} - x = 4$

 $x = \frac{4}{\sqrt{5} - 1}$

 $x = \frac{4(\sqrt{5} + 1)}{5 - 1^2}$

 $x = 1 + \sqrt{5}$

 The desired quantity is $1 + 5 = 6$, and the answer is $\boxed{A}$.

 


	\item (\textit{AHSME 1997}) A list of integers has mode $32$ and mean $22$. The smallest number in the list is $10$. The median $m$ of the list is a member of the list. If the list member $m$ were replaced by $m+10$, the mean and median of the new list would be $24$ and $m+10$, respectively. If were $m$ instead replaced by $m-8$, the median of the new list would be $m-4$. What is $m$?


 $\textbf{(A)}\ 16\qquad\textbf{(B)}\ 17\qquad\textbf{(C)}\ 18\qquad\textbf{(D)}\ 19\qquad\textbf{(E)}\ 20$


 
 \textbf{Answer: }E



\textbf{Solution}

Let there be $n$ integers on the list.  The list of $n$ integers has mean $22$, so the sum of the integers is $22n$.

 Replacing $m$ with $m+10$ will increase the sum of the list from $22n$ to $22n + 10$.

 The new mean of the list is $24$, so the new sum of the list is also $24n$.

 Thus, we get $22n + 10 = 24n$, leading to $n=5$ numbers on the list.

 If there are $5$ numbers on the list with mode $32$ and smallest number $10$, then the list is $\{10, x, m, 32, 32\}$

 Since replacing $m$ with $m-8$ gives a new median of $m-4$, and $m-4$ must be on the list of $5$ integers since $5$ is odd, $x = m-4$, and the list is now $\{10, m-4, m, 32, 32\}$

 The sum of the numbers on this list is $22n = 22\cdot 5 = 110$, so we get:

 $10 + m - 4 + m + 32 + 32 = 110$

 $70 + 2m = 110$

 $m = 20$, giving answer $\boxed{E}$.

 The original list is $\{10, 16, 20, 32, 32\}$, with mean $\frac{10 + 16 + 20 + 32 + 32}{5} = 22$ and median $20$ and mode $32$.

 The second list is $\{10, 16, 30, 32, 32\}$, with mean $\frac{10 + 16 + 30 + 32 + 32}{5} = 24$ and median $m + 10 = 30$.

 The third list is $\{10, 16, 12, 32, 32\} \rightarrow \{10, 12, 16, 32,32\}$ with median $m-4 = 16$.

 


	\item (\textit{AHSME 1997}) A circle with center $O$ is tangent to the coordinate axes and to the hypotenuse of the $30^\circ$-$60^\circ$-$90^\circ$ triangle $ABC$ as shown, where $AB=1$. To the nearest hundredth, what is the radius of the circle?


 


                    \begin{figure}[H]
                        centering
                        \includegraphics[width=0.5\textwidth]{Problem/AHSME_I_1997_19_Problem_Figure_1.png}
                    \end{figure}
                    


 $\textbf{(A)}\ 2.18\qquad\textbf{(B)}\ 2.24\qquad\textbf{(C)}\ 2.31\qquad\textbf{(D)}\ 2.37\qquad\textbf{(E)}\ 2.41$


 
 \textbf{Answer: }D



\textbf{Solution}


                    \begin{figure}[H]
                        centering
                        \includegraphics[width=0.5\textwidth]{Solution/AHSME_I_1997_19_Solution_Figure_1.png}
                    \end{figure}
                    

 Draw radii $OE$ and $OD$ to the axes, and label the point of tangency to triangle $ABC$ point $F$.  Let the radius of the circle $O$ be $r$.  Square $OEAD$ has side length $r$.

 Because $BD$ and $BF$ are tangents from a common point $B$, $BD = BF$.

 $AD = AB + BD$

 $r = 1 + BD$

 $r = 1 + BF$

 Similarly, $CF = CE$, and we can write:

 $AE = AC + CE$

 $r = \sqrt{3} + CF$

 Equating the radii lengths, we have $1 + BF = \sqrt{3} + CF$

 This means $BF - CF = \sqrt{3} - 1$

 $BF + CF = 2$ by the 30-60-90 triangle.

 Therefore, $2BF = 2 + \sqrt{3} - 1$, and we get $BF = \frac{1}{2} + \frac{\sqrt{3}}{2}$

 The radius of the circle is $AD$, which is $BF + 1 = \frac{3}{2} + \frac{\sqrt{3}}{2}$

 Using decimal approximations, $r \approx 1.5 + \frac{1.73^+}{2} \approx 2.37$, and the answer is $\boxed{D}$.

 



\textbf{Solution 2}


From the diagram above, it is more direct to note that BC = CF + BF = r - $\sqrt{3}$ + r - 1 = 2

 



\textbf{Solution 3}

The total area of both kites is $BC \cdot OF= 2r$. Thus, $r^2 = \frac{\sqrt{3}}{2} + 2r$

 


	\item (\textit{AHSME 1997}) Which one of the following integers can be expressed as the sum of $100$ consecutive positive integers?


 $\textbf{(A)}\ 1,\!627,\!384,\!950\qquad\textbf{(B)}\ 2,\!345,\!678,\!910\qquad\textbf{(C)}\ 3,\!579,\!111,\!300\qquad\textbf{(D)}\ 4,\!692,\!581,\!470\qquad\textbf{(E)}\ 5,\!815,\!937,\!260$


 
 \textbf{Answer: }A



\textbf{Solution}

The sum of the first $100$ integers is $\frac{100\cdot 101}{2} = 5050$.

 If you add an integer $k$ to each of the $100$ numbers, you get $5050 + 100k$, which is the sum of the numbers from $k+1$ to $k+100$.

 You're only adding multiples of $100$, so the last two digits will remain unchanged.

 Thus, the only possible answer is $\boxed{A}$, because the last two digits are $50$.

 As an aside, if $5050 + 100k = 1627384950$, then $k = 16273799$, and the numbers added are the integers from $16273800$ to $16273899$.

 



\textbf{Solution}

Notice how the sum of 100 consecutive integers is $(x-49)+(x-48)+(x-47)...+x+...(x+47)+(x+48)+(x+49)+(x+50)$. 

 Cancelling out the constants give us $100x + 50$. 

 Looking over at the list of possible values, we quickly realise that the only possible solution is $\boxed{A}$

 


	\item (\textit{AHSME 1997}) For any positive integer $n$, let 


 $f(n) =\left\{\begin{matrix}\log_{8}{n}, &\text{if }\log_{8}{n}\text{ is rational,}\\ 0, &\text{otherwise.}\end{matrix}\right.$


 What is $\sum_{n = 1}^{1997}{f(n)}$?


 $\textbf{(A)}\ \log_{8}{2047}\qquad\textbf{(B)}\ 6\qquad\textbf{(C)}\ \frac{55}{3}\qquad\textbf{(D)}\ \frac{58}{3}\qquad\textbf{(E)}\ 585$


 
 \textbf{Answer: }C



\textbf{Solution}

For positive integers $n$, $\log_8 n$ is rational if and only if $n = 2^k$ for an integer $k$.  That's because $\log_8 2^k = k\log_8 2 = \frac{k}{3}$.

 So we actually want to find $\sum_{k=0}^{10} \log_8 2^k$, since $2^{11}$ will be over $1997$.

 Using log properties, we get $\sum_{k=0}^{10} k \log_8 2$

 $\frac{1}{3}\sum_{k=0}^{10} k$

 $\frac{1}{3}\cdot (\frac{10\cdot 11}{2})$

 $\frac{55}{3}$, and the answer is $\boxed{C}$

 


	\item (\textit{AHSME 1997}) Ashley, Betty, Carlos, Dick, and Elgin went shopping. Each had a whole number of dollars to spend, and together they had $56$ dollars. The absolute difference between the amounts Ashley and Betty had to spend was $19$ dollars. The absolute difference between the amounts Betty and Carlos had was $7$ dollars, between Carlos and Dick was $5$ dollars, between Dick and Elgin was $4$ dollars, and between Elgin and Ashley was $11$ dollars. How many dollars did Elgin have?


 $\textbf{(A)}\ 6\qquad\textbf{(B)}\ 7\qquad\textbf{(C)}\ 8\qquad\textbf{(D)}\ 9\qquad\textbf{(E)}\ 10$


 
 \textbf{Answer: }E



\textbf{Solution}

Working backwards, if $6 \le E \le 10$, then $6 \pm 11 \le A \le 10 \pm 11$.  Since $A$ is a positive integer, $17 \le A \le 21$. 

 
Since $17 \le A \le 21$, we know that $17 \pm 19 \le B \le 21 \pm 19$.  But if $B=36$, which is the smallest possible "plus" value, then $E + A + B = 6 + 17 + 36 = 59$, which is too much money.

 
Hence, $17 - 19 \le B \le 21 - 19$.  But since $B$ must be a positive integer, that leaves only two possibilities:  $B = 1$ or $B=2$, which correspond with $E = 9$ and $E = 10$.

 
Concentrating only on $E=9$, we have $E=9$ leading to $A = 9 + 11 = 20$, which leads to $B = 20 - 19 = 1$, which leads to $C = 1 + 7 = 8$.  Thus far we have given out $9 + 20 + 1 + 8 = 38$ dollars.  This means that Dick must have $56 - 38 = 18$ dollars.  However, the difference between Carlos and Dick is not $5$ dollars.

 
Thus, the right answer must be $\boxed{\textbf{(E)}}$.  Verifying, if $E = 10$, then $A = 10 + 11 = 21$, $B = 21 - 19 = 2$, which leads to $C = 2 + 7 = 9$.  Thus far, we have given out $10 + 21 + 2 + 9 = 42$ dollars, leaving $56 - 42 = 14$ dollars for Dick.  Dick does indeed have $5$ dollars more than Carlos, and $4$ dollars more than Elgin.

 


	\item (\textit{AHSME 1997}) 
                    \begin{figure}[H]
                        centering
                        \includegraphics[width=0.5\textwidth]{Problem/AHSME_I_1997_23_Problem_Figure_1.png}
                    \end{figure}
                    


 In the figure, polygons $A$, $E$, and $F$ are isosceles right triangles; $B$, $C$, and $D$ are squares with sides of length $1$; and $G$ is an equilateral triangle. The figure can be folded along its edges to form a polyhedron having the polygons as faces. The volume of this polyhedron is


 $\textbf{(A)}\ 1/2\qquad\textbf{(B)}\ 2/3\qquad\textbf{(C)}\ 3/4\qquad\textbf{(D)}\ 5/6\qquad\textbf{(E)}\ 4/3$


 
 \textbf{Answer: }D



\textbf{Solution}

Since there are three squares, the figure may appear to be based upon a unit cube.

 Since there is an equilateral triangle of side length $\sqrt{2}$, this triangle has the length of a facial (non-spatial) diagonal of a unit cube.

 Since there are three isosceles right triangles that are half the area of a face, they might represent faces of the cube that have been cut in half along the diagonal.

 Putting these three clues together, we can mentally fold the solid back together to form a unit cube that has had one vertex removed.  Along with this vertex, the triangular half of each of the three adjoining faces has been removed.  When the unit cube is sliced like this, it will cut three of the faces into isosceles triangles like $A$, $E$, and $F$.  Additionally, it will create an equilateral triangle that has as its sides three facial diagonals, just like $G$.

 So the volume is a unit cube with a tetrahedron removed.  The tetrahedron has three right triangular faces like $A$, $E$, and $F$, and one equilateral triangular face like $G$.  The volume of this removed tetrahedron can be calculated as follows:

 Place one of the three right triangles on the ground as the base.  It has an area of $\frac{1}{2}$.  The height is the height of one of the other two right triangles - the equilateral triangle is slanted and is not a height.  Thus, the volume of the tetrahedron is $\frac{1}{3}\cdot \frac{1}{2}\cdot 1 = \frac{1}{6}$, and when this volume is removed, the remaining volume is $\frac{5}{6}$, which is option $\boxed{D}$.

 


	\item (\textit{AHSME 1997}) A rising number, such as $34689$, is a positive integer each digit of which is larger than each of the digits to its left. There are $\binom{9}{5} = 126$ five-digit rising numbers. When these numbers are arranged from smallest to largest, the $97^{th}$ number in the list does not contain the digit


 $\textbf{(A)}\ 4\qquad\textbf{(B)}\ 5\qquad\textbf{(C)}\ 6\qquad\textbf{(D)}\ 7\qquad\textbf{(E)}\ 8$


 
 \textbf{Answer: }B



\textbf{Solution}

The list starts with $12345$.  There are $\binom{8}{4} = 70$ four-digit rising numbers that do not begin with $1$, and thus also $70$ five digit rising numbers that do begin with $1$ that are formed by simply putting a $1$ before the four digit number.

 Thus, the $71^{\text{st}}$ number is $23456$.  There are $\binom{6}{3} = 20$ three-digit rising numbers that do not begin with a $1,2$ or $3$, and thus $20$ five digit rising numbers that begin with a $23$.

 Thus, the $91^{\text{st}}$ number is $24567$.  Counting up, $24568, 24569, 24578, 24579, 24589, 24678$ is the $97^{\text{th}}$ number, which does not contain the digit $\boxed{\textbf{(B)} \ 5}$.

 


	\item (\textit{AHSME 1997}) Let $ABCD$ be a parallelogram and let $\overrightarrow{AA^\prime}$, $\overrightarrow{BB^\prime}$, $\overrightarrow{CC^\prime}$, and $\overrightarrow{DD^\prime}$ be parallel rays in space on the same side of the plane determined by $ABCD$. If $AA^\prime = 10$, $BB^\prime = 8$, $CC^\prime = 18$, and $DD^\prime = 22$ and $M$ and $N$ are the midpoints of $A^\prime C^\prime$ and $B^\prime D^\prime$, respectively, then $MN =$


 $\textbf{(A)}\ 0\qquad\textbf{(B)}\ 1\qquad\textbf{(C)}\ 2\qquad\textbf{(D)}\ 3\qquad\textbf{(E)}\ 4$


 
 \textbf{Answer: }B



\textbf{Solution}

Let $ABCD$ be a unit square with $A(0,0,0)$, $B(0,1,0)$, $C(1,1,0)$, and $D(1,0,0)$.  Assume that the rays go in the +z direction.  In this case, $A^\prime(0,0,10)$, $B^\prime(0,1,8)$,  $C^\prime(1,1,18)$, and $D^\prime(1,0,22)$.  Finding the midpoints of $A^\prime C^\prime$ and $B^\prime D^\prime$ gives $M(\frac{1}{2}, \frac{1}{2}, 14)$ and $N(\frac{1}{2}, \frac{1}{2}, 15)$.  The distance $MN$ is $15 - 14 = 1$, and the answer is $\boxed{B}$.    

 
While this solution is not a proof, the problem asks for a value $MN$ that is supposed to be an invariant.  Additional assumptions that are consistient with the problem's premise should not change the answer.  The assumptions made in this solution are that $ABCD$ is a square, and that all rays $XX^\prime$ are perpendicular to the plane of the square.  These assumptions are consistient with all facts of the problem.  Additionally, the assumptions allow for a quick solution to the problem.

 


	\item (\textit{AHSME 1997}) Triangle $ABC$ and point $P$ in the same plane are given. Point $P$ is equidistant from $A$ and $B$, angle $APB$ is twice angle $ACB$, and $\overline{AC}$ intersects $\overline{BP}$ at point $D$. If $PB = 3$ and $PD= 2$, then $AD\cdot CD =$


 
                    \begin{figure}[H]
                        centering
                        \includegraphics[width=0.5\textwidth]{Problem/AHSME_I_1997_26_Problem_Figure_1.png}
                    \end{figure}
                    


 $\textbf{(A)}\ 5\qquad\textbf{(B)}\ 6\qquad\textbf{(C)}\ 7\qquad\textbf{(D)}\ 8\qquad\textbf{(E)}\ 9$


 
 \textbf{Answer: }A



\textbf{Solution}

The product of two lengths with a common point brings to mind the \href{https://artofproblemsolving.com/wiki/index.php?title=Power\_of\_a\_Point\_Theorem}{Power of a Point Theorem}.

 Since $PA = PB$, we can make a circle with radius $PA$ that is centered on $P$, and both $A$ and $B$ will be on that circle.  Since $\angle APB = \widehat {AB} = 2 \angle ACB$, we can see that point $C$ will also lie on the circle, since the measure of arc $\widehat {AB}$ is twice the measure of inscribed angle $\angle ACB$, which is true for all inscribed angles.

 
                    \begin{figure}[H]
                        centering
                        \includegraphics[width=0.5\textwidth]{Solution/AHSME_I_1997_26_Solution_Figure_1.png}
                    \end{figure}
                    

 Since $PDB$ is a line, we have $PD + DB = PB$, which gives $3 = DB + 2$, or $DB = 1$.

 We now extend radius $PB = 3$ to diameter $EB = 6$.  Since $EDB$ is a line, we have $ED + DB = EB$, which gives $ED + 1 = 6$, or $ED = 5$.

 Finally, we apply the power of a point theorem to point $D$.  This states that $AD \cdot DC = DB \cdot DE$.  Since $DB = 1$ and $DE = 5$, the desired product is $5$, which is $\boxed{A}$.

 



\textbf{Solution 2}

Construct the angle bisector of $\angle APD,$ and let it intersect $AD$ at $E.$ From the angle bisector theorem, we have $AE=3a$ and $DE=2a$ for some $a.$ Then, note that $\angle EPD = \angle BCD = x,$ so $EPCB$ is cyclic. Then, $\frac{PD}{ED} = \frac{CD}{BD}$ or $\frac{2}{2x} = \frac{CD}{1}.$ Thus, $AD \cdot DC = 5x \cdot DC = 5,$ or $\boxed{A}.$

 


	\item (\textit{AHSME 1997}) Consider those functions $f$ that satisfy $f(x+4)+f(x-4) = f(x)$ for all real $x$. Any such function is periodic, and there is a least common positive period $p$ for all of them. Find $p$.


 $\textbf{(A)}\ 8\qquad\textbf{(B)}\ 12\qquad\textbf{(C)}\ 16\qquad\textbf{(D)}\ 24\qquad\textbf{(E)}\ 32$


 
 \textbf{Answer: }D



\textbf{Solution}

Recall that $p$ is the fundamental period of function $f$ iff $p$ is the smallest positive $p$ such that $f(x) = f(x + p)$ for all $x$.

 
In this case, we know that $f(x+ 4) + f(x - 4) = f(x)$.  Plugging in $x+4$ in for $x$ to get the next equation in the recursion, we also get $f(x + 8) + f(x) = f(x + 4)$.  Adding those two equations gives $f(x + 8) +  f(x - 4) = 0$ after cancelling out common terms.

 
Again plugging in $x + 4$ in for $x$ in that last equation (in order to get $f(x)$), we find that $f(x) = -f(x + 12)$.  Now, plugging in $x+12$ for $x$, we get $f(x + 12) = -f(x + 24)$.  This proves that $f(x)  = f(x + 24)$, so there is a period of $24$, which gives answer $\boxed{D}$.  We now eliminate answers $A$ through $C$.

 
Let $f(x) = x$ for $x \in \{1, 2, 3, 4, 5, 6, 7, 8\}$.  Plugging in $x=5$ into the initial equation gives $f(9) + f(1) = f(5)$, which implies that $f(9) = 4$.  Since $f(1) \neq f(1 + 8)$, the function does not have period $8$.

 Continuing, $f(10) = f(6) - f(2) = 4$, $f(11) = f(7) - f(3) = 4$, and $f(12) = f(8) - f(4) = 4$.  Once we hit $f(13)$, we have $f(13) = f(9) - f(5) = 4 - 5 = -1$.  Since $f(1) \neq f(1 + 12)$, the function does not have period $12$.

 Finally, $f(14) = f(10) - f(6) = 4 - 6 = -2$, $f(15) = f(11) - f(7) = 4 - 7 = -3$, $f(16) = f(12) - f(8) = 4 - 8 = -4$, and $f(17) = f(13) - f(9) = -1 - 4 = -5$.  Since $f(1) \neq f(1 + 16)$, the function does not have period $16$.

 To confirm that our original period works, we may see that $f(18) = -6$, $f(19) = -7$, $f(20) = -8$, $f(21) = -4$, $f(22) = -4$, $f(23) = -4$, and $f(24) = -4$.  Finally, $f(25) = 1$, which is indeed the same as $f(1)$.

 


	\item (\textit{AHSME 1997}) How many ordered triples of integers $(a,b,c)$ satisfy $|a+b|+c = 19$ and $ab+|c| = 97$? 


 $\textbf{(A)}\ 0\qquad\textbf{(B)}\ 4\qquad\textbf{(C)}\ 6\qquad\textbf{(D)}\ 10\qquad\textbf{(E)}\ 12$


 
 \textbf{Answer: }E



\textbf{Solution}

\href{https://artofproblemsolving.com/wiki/index.php?title=WLOG}{WLOG}, let $a \ge 0$, and let $a \ge b$.  We can say this because if we have one solution $(a,b) = (a_0, b_0)$ with $a_0 \ge 0$ and $a_0 > b_0$, we really have the four solutions $(a_0, b_0), (-a_0, -b_0), (b_0, a_0), (-b_0, -a_0)$ by the symmetry of the original problem.  

 Furthermore, we assert that these four solutions are distinct.  We can say that $a > b$, since if $a=b$, we have $c = 19 - 2a$ for the first equation and either $c = 97 - a^2$ or $c = a^2 - 97$ for the second equation.  Equating $19 - 2a = 97 - a^2$ gives no integer solution, while equating $19 - 2a = a^2 - 97$ also gives no integer solution.

 Thus, we can now assume WLOG that $a \ge 0$ and $a > b$, and each pair of $(a_0,b_0)$ that we get will generate four unique solutions:  $(a_0, b_0), (b_0, a_0), (-a_0, -b_0), (-b_0, -a_0)$.

 
We now divide the problem into $c \ge 0$ and $c < 0$:

 If $c \ge 0$, we have $|a + b| + c = 19$ and $ab + c = 97$.

 Solving both equations for $c$ and equating them, we get that $ab - |a + b| = 78$.  Splitting these up, we find that either $ab - a - b = 78$ or $ab + a+ b = 78$.  Factoring both with \href{https://artofproblemsolving.com/wiki/index.php?title=SFFT}{SFFT} gives $(a-1)(b-1) = 79$ or $(a +1)(b+1) = 79$.  We factor with the restrctions that $a \ge 0$ and $a > b$.  Since $79$ is prime, we have:

 $a - 1 = 79$ and $b - 1 = 1$, which leads to $(80,2)$.

 $a + 1 = 79$ and $b + 1 = 1$, which leads to $(78, 0)$.

 Each of those solutions could generate $3$ more solutions, giving a total of $8$ potential solutions.  However, in each the first set of four solutions, we have $|a + b| = 82$, which from the original first equation $|a + b| + c = 19$ gives $c = 19 - 82$, which contradicts our initial assumption that $c\ge 0$.  Similarly, for the second set of four solutions, we have $|a + b| = 78$, which leads to $c = 19 - 78$, also contradicting $c \ge 0$.

 
If $c < 0$, we have $|a + b| + c = 19$ and $ab - c = 97$.  We note that $19 - c$ must be positive whenever $c$ is negative, and thus $|a + b| = a + b$.

 Solving both equations for $c$ and using SFFT as above gives $(a+1)(b+1) = 117$.  Since $117 = 3^2\cdot 13$, we factor with the restriction that $a\ge 0$ and $a > b$.  Thus, we can let $a+1  \in \{117, 39, 13\}$, which means $a \in \{116, 38, 12\}$.  These give corresponding $b+1\in \{1, 3, 9\}$, which leads to corresponding $b \in \{0, 2, 8\}$.  Combining the solutions, we have $(a,b) = (116, 0), (38, 2), (12, 8)$.  

 Each of these three solutions permutes, negates, and permute-negates into $4$ solutions as described in the start of the solution, for a total of $12$ solutions.

 Checking our solutions to ensure $c < 0$, we find in the first set of four solutions, $|a + b| = 116$, and thus $c = 19 - 116$, which is indeed negative.

 In the second set of four solutions, $|a + b| = 40$, which leads to $c = 19 - 40$, which is also negative.

 Finally, in the third set of four solutoins, $|a + b| = 20$, which leads to $c = 19 - 20$, which is negative.

 Thus, there are $12$ ordered triples, and the answer is $\boxed{E}$.

 


	\item (\textit{AHSME 1997}) Call a positive real number special if it has a decimal representation that consists entirely of digits $0$ and $7$. For example, $\frac{700}{99}= 7.\overline{07}= 7.070707\cdots$ and $77.007$ are special numbers. What is the smallest $n$ such that $1$ can be written as a sum of $n$ special numbers?


 $\textbf{(A)}\ 7\qquad\textbf{(B)}\ 8\qquad\textbf{(C)}\ 9\qquad\textbf{(D)}\ 10\qquad\\ \textbf{(E)}\ \text{The number 1 cannot be represented as a sum of finitely many special numbers.}$


 
 \textbf{Answer: }B



\textbf{Solution}

Define a super-special number to be a number whose decimal expansion only consists of $0$'s and $1$'s. The problem is equivalent to finding the number of super-special numbers necessary to add up to $\frac{1}{7}=0.142857142857\hdots$. This can be done in $8$ numbers if we take
 \[0.111111\hdots, 0.011111\hdots, 0.010111\hdots, 0.010111\hdots, 0.000111\hdots, 0.000101\hdots, 0.000101\hdots, 0.000100\hdots\]Now assume for sake of contradiction that we can do this with strictly less than $8$ super-special numbers (in particular, less than $10$.) Then the result of the addition won't have any carry over, so each digit is simply the number of super-special numbers which had a $1$ in that place. This means that in order to obtain the $8$ in $0.1428\hdots$, there must be $8$ super-special numbers, so the answer is $\boxed{\textbf{(B)}\ 8}$.

 


	\item (\textit{AHSME 1997}) For positive integers $n$, denote $D(n)$ by the number of pairs of different adjacent digits in the binary (base two) representation of $n$. For example, $D(3) = D(11_{2}) = 0$, $D(21) = D(10101_{2}) = 4$, and $D(97) = D(1100001_{2}) = 2$. For how many positive integers less than or equal to $97$ does $D(n) = 2$?


 $\textbf{(A)}\ 16\qquad\textbf{(B)}\ 20\qquad\textbf{(C)}\ 26\qquad\textbf{(D)}\ 30\qquad\textbf{(E)}\ 35$


 
 \textbf{Answer: }C



\textbf{Solution 1}

If $D(n)$ is even, then the binary expansion of $n$ will both begin and end with a $1$, because all positive binary numbers begin with a $1$, and if you switch digits twice, you will have a $1$ at the end.  Thus, we are only concerned with the $49$ odd numbers between $1$ and $98$ inclusive.

 All of these odd numbers will have an even $D(n)$.  $D(n) = 0$ will be given by the numbers $1, 11, 111, 1111, 11111, 111111$, which is a total of $6$ numbers.

 We skip $D(n) = 2$ for now, and move to $D(n) = 4$, which is easier to count.  The smallest $D(n) = 4$ happens when $n = 10101$.  To get another number such that $D(n) = 4$, we may extend any of the five blocks of zeros or ones by one digit.  This will form $110101, 100101, 101101, 101001, 101011$, all of which are odd numbers that have $D(n) = 4$.  To find seven digit numbers that have $D(n) = 4$, we can again extend any block by one, so long as it remains less than $1100001$ or under.  There are five cases.

 1)  Extending $110101$ is impossible without going over $1100001$.

 2)  Extending $100101$ by putting a $1$ at the beginning will go over $1100001$, but the other four extensions work, giving $1000101, 1001101, 1001001, 1001011$.

 3)  Extending $101101$ by putting a $1$ at the beginning will go over $1100001$, but the other four extensions give $1001101, 1011101, 1011001, 1011011$.  However, $1001101$ already appeared in \#2, giving only three new numbers.

 4)  Extending $101001$ at the first group is impossible.  The other four extensions are $1001001, 1011001, 1010001, 1010011$, but the first two are repeats.  Thus, there are only two new numbers.

 5)  Extending $101011$ at the first group is impossible.  The other four extensions give $1001011, 1011011, 1010011, 1010111$, but only the last number is new.

 Thus, there is $1$ five digit number, $5$ six digit numbers, and $4 + 3 + 2 + 1 = 10$ seven digit numbers under $1100001$ for which $D(n) = 4$.  That gives a total of $16$ numbers.

 There smallest number for which $D(n) = 6$ is $1010101$, which is under $98$.  Further extensions, as well as cases where $D(n) > 6$, are not possible.

 Thus, we know that there are $6$ odd numbers that have $D(n) = 0$, and $16$ odd numbers that have $D(n) = 4$, and $1$ number that has $D(n) = 6$.  The remaining odd numbers must have $D(n) = 2$.  This means there are $49 - 6 - 16 - 1 = 26$ numbers that have $D(n) = 2$, which is option $\boxed{C}$

 



\textbf{Solution 2}

Each pair of different adjacent digits corresponds to a change from one a chain of 1s to a chain of 0s or vice versa. For example, in the number $111001$ there are 2 such pairs and also 2 changes (three 1s change to two 0s, then change back to 1). Thus, in a number with 2 changes, there will be three "blocks" of consecutive digits. Since the number always starts with a 1, each valid number will have the form 
 \[1\dots0\dots1\dots\]

 The total number of digits cannot exceed $7$. If the number has 6 digits in binary or less: this is equivalent to solving $a+b+c \le 6$, where $a, b, c$ represent how many 1s, 0s, and 1s are in the three blocks where $a, b, c$ are positive integers. Using stars and bars, we find that there are $\binom{6}{3}=20$ such binary numbers.

 If the number has $7$ digits: we see that there are $5$ numbers for which the second digit is 0, and only 1 valid number for which the second digit is 1 (97 is 1100001 in binary). Thus, we have a total of $20+5+1=\fbox{(C) 26}$ numbers.

 



\textbf{Solution 3 (Casework)}

For $D(n)$ to be $2$, $n$ must be in the form $1...10...01...1$. Note that $n$ must be at least $3$ digits for $D(n)$ to be $2$

 Case $1$: $n$ has $3$ digits

 $n$ = $101_2$, there is only $1$ possible value for $n$ when $n$ has $3$ digits.

 
Case $2$: $n$ has $4$ digits

 $n = 1$ \_ \_ $1_2$, there could be $1$ zero or $2$ zeros between the first digit and the last digit. If there is only $1$ zero, the zero has $2$ possible places. If there are $2$ zeros, there is only $1$ possible placement. Therefore, there are $2+1=3$ possible values for $n$ when $n$ has $4$ digits.

 
Case $3$: $n$ has $5$ digits

 $n = 1$ \_ \_ \_ $1_2$, there could be $1$, $2$ or $3$ zeros between the first digit and the last digit. If there is only $1$ zero, the zero has $3$ possible places. If there are $2$ zeros, there are $2$ possible placements. If there are $3$ zeros, there is only $1$ possible placement. Therefore, there are $3+2+1=6$ possible values for $n$ when $n$ has $5$ digits.

 
Case $4$: $n$ has $6$ digits

 $n = 1$ \_ \_ \_ \_ $1_2$, there could be $1$, $2$, $3$ or $4$ zeros between the first digit and the last digit. If there is only $1$ zero, the zero has $4$ possible places. If there are $2$ zeros, there are $3$ possible placements. If there are $3$ zeros, there are $2$ possible placements. If there are $4$ zeros, there is only $1$ possible placement. Therefore, there are $4+3+2+1=10$ possible values for $n$ when $n$ has $6$ digits.

 
Case $4$: $n$ has $7$ digits

 Given by the problem, $n \le 1100001_2$. We could first count the number of $n$ that is in the form of $10$ \_ \_ \_ \_ $1_2$, where the digits between the second and last digit are in the form of $0...01...1$. 

 There could be $0$, $1$, $2$, $3$ or $4$ zeros. Therefore, there are $5$ possible values for $n$ when $n=10$ \_ \_ \_ \_ $1_2$.

 There is also the case were $n = 1100001_2$. Therefore, there are $5+1=6$ possible values for $n$ when $n$ has $7$ digits.

 Therefore, the answer is $1+3+6+10+6=\boxed{\textbf{(C) } 26}$.

 small edit ~\href{https://artofproblemsolving.com/wiki/index.php/User:Mathlete0001}{Mathlete0001}

 \textbf{Sidenote from samrocksnature} Note that any binary number $n_2$ must start with $1,$ end with $1,$ and have some number of $0$s in between in order for $D(n)=2.$ Let $n_2$ have $a$ leading $1$s, $b$ middle $0$s, and $c$ trailing $1$s. Also let $n_2$ have $k$ digits. 

 Lemma 1: There are $\binom{k-1}{2}$ positive integers $n_2$ with $k$ digits such that $D(n)=2.$ 

 There are $k-2$ slots to fill in with $a+b+c$ digits, where $a,c\geq 0$ and $b \geq 1.$ Let $b'=b-1$ such that $b \geq 0,$ since we must have at least $0$ in our $k-2$ slots. Thus, there are now $k-1$ slots to fill in with $a+b'+c$ digits for any nonnegative integers $a,b',c.$ By stars and bars, there are $\binom{k-1}{2}$ ways to do this. 

 It makes sense that $k \geq 3,$ since $101_2$ is the smallest such $n_2$ that satisfies $D(n)=2.$

 Lemma 2: There are $\frac{x(x-1)(x-2)}{6}$ ways for $D(n)=2$ when $n \leq 2^{x+1}-1$ for some nonnegative integer $x.$

 This inequality condition for $n$ translates to $n$ having less than or equal to $x$ digits in base $2.$ By Lemma 1, the sum we seek is $\sum_{k=3}^{x}\binom{k-1}{2}=1+3+6+ \cdots + \text{ (x-2)th triangular number}.$ By telescoping or prior knowledge of the formula, we obtain our lemma.

 


\end{enumerate}